\chapter{System Implementation and Experimental Evaluation}
\label{chap:system-architecture}

This chapter describes the implemented monitoring system, the experimental setup, and the results obtained on the fixed monitoring subset in the observed time window.

\section{Overview}
\label{sec:evaluation-pipeline-overview}

The system has five layers: Langfuse datasets (\(\mathcal{D}_N\), \(\mathcal{D}_n\)); orchestration via \texttt{run\_eval.py} into one canonical request; provider adapters (OpenAI, Gemini, Claude); retrieval/ingestion into Langfuse traces; export/dashboard CSVs. Three GitHub Actions workflows implement the weekly loop (Figure~\ref{fig:ga-weekly-pipelines}): submit batches, retrieve and ingest, sync the Streamlit dashboard.

Routine automation uses only the frozen \(n=300\) panel (\(\approx\) USD~20/pass). Full-benchmark passes (\(N=1199\), \(\approx\) USD~90) are for baseline initialization and explicit escalation. The roster and prices live in \texttt{config/models.yaml} (Table~\ref{tab:monitored-models}).

\begin{table}[htbp]
\centering
\caption{LLM models monitored on the evaluation pipelines, grouped by provider.}
\label{tab:monitored-models}
\begin{tabular}{ll}
\toprule
Provider & Models \\
\midrule
OpenAI   & \texttt{gpt-5.4-2026-03-05}, \texttt{gpt-5.2}, \texttt{gpt-5.1}, \texttt{gpt-5}, \\
         & \texttt{gpt-5-mini}, \texttt{gpt-5-nano}, \texttt{gpt-4.1-mini}, \texttt{gpt-4.1-nano}. \\
Gemini   & \texttt{gemini-2.5-flash}. \\
Claude   & \texttt{claude-sonnet-4-6}, \texttt{claude-sonnet-4-5}, \texttt{claude-haiku-4-5}. \\
\bottomrule
\end{tabular}
\end{table}

\section{Implemented Monitoring Pipeline}

Provider batch APIs are asynchronous, so submission and retrieval are separate jobs. Figure~\ref{fig:ga-weekly-pipelines} shows the weekly schedule: \texttt{weekly\_drift\_monitoring.yml} (Fri 05:00 UTC), \texttt{weekly\_retrieve.yml} (Mon 05:00 UTC), \texttt{sync-streamlit-dashboard.yml} (Mon 07:00 UTC).

\begin{figure}[htbp]
    \centering
    \begin{adjustbox}{max width=\linewidth}
    % Scheduled GitHub Actions workflows. Each lane is one workflow; the schedule is
% carried in the lane label rather than in a node.
\begin{tikzpicture}[
  font=\scriptsize,
  >=Latex,
  wfstep/.style={
    draw=vizaxis, line width=0.4pt, rounded corners=1.5pt, fill=white,
    align=center, text width=21mm, minimum height=10mm, inner sep=2pt
  },
  lane/.style={draw=vizgrid, line width=0.5pt, rounded corners=2pt},
  lanelabel/.style={font=\scriptsize\bfseries, text=vizink, anchor=west, fill=white, inner xsep=1.5mm, inner ysep=0.4mm},
  flow/.style={->, vizmuted, line width=0.5pt},
  hand/.style={->, vizblue, line width=0.6pt},
  node distance=3.8mm,
]

% ---------------- lane 1: submit
\node[wfstep] (a1) {Run evaluation\\\texttt{run\_eval.py}};
\node[wfstep, right=of a1] (a2) {Load panel\\from Langfuse};
\node[wfstep, right=of a2] (a3) {Serialize through\\provider adapters};
\node[wfstep, right=of a3] (a4) {Submit provider\\batch jobs};
\node[wfstep, right=of a4] (a5) {Cache batch\\manifest};
\foreach \i/\j in {a1/a2, a2/a3, a3/a4, a4/a5} \draw[flow] (\i) -- (\j);

% ---------------- lane 2: retrieve
\node[wfstep, below=18mm of a1] (b1) {Restore cached\\batch manifest};
\node[wfstep, right=of b1] (b2) {Poll batch jobs\\and retrieve results};
\node[wfstep, right=of b2] (b3) {Persist raw results\\as JSONL};
\node[wfstep, right=of b3] (b4) {Join labels and\\compute scores};
\node[wfstep, right=of b4] (b5) {Store traces and\\scores in Langfuse};
\foreach \i/\j in {b1/b2, b2/b3, b3/b4, b4/b5} \draw[flow] (\i) -- (\j);

% ---------------- lane 3: dashboard
\node[wfstep, below=18mm of b1] (c1) {Export Langfuse\\CSV snapshot};
\node[wfstep, right=of c1] (c2) {Commit versioned\\analysis snapshot};
\node[wfstep, right=of c2] (c3) {Publish data to\\dashboard repository};
\foreach \i/\j in {c1/c2, c2/c3} \draw[flow] (\i) -- (\j);

% ---------------- lane frames and labels
\begin{scope}[on background layer]
  \node[lane, fit=(a1)(a5), inner xsep=3mm, inner ysep=3.2mm] (L1) {};
  \node[lane, fit=(b1)(b5), inner xsep=3mm, inner ysep=3.2mm] (L2) {};
  \node[lane, fit=(c1)(c3)(b5|-c1), inner xsep=3mm, inner ysep=3.2mm] (L3) {};
\end{scope}
\node[lanelabel] at ($(L1.north west)+(2mm,0)$)
  {\texttt{weekly\_drift\_monitoring.yml} \textnormal{--- Fri 05:00 UTC}};
\node[lanelabel] at ($(L2.north west)+(2mm,0)$)
  {\texttt{weekly\_retrieve.yml} \textnormal{--- Mon 05:00 UTC}};
\node[lanelabel] at ($(L3.north west)+(2mm,0)$)
  {\texttt{sync-streamlit-dashboard.yml} \textnormal{--- Mon 07:00 UTC}};

% ---------------- cross-lane handoffs, routed through the inter-lane gap
\coordinate (mid1) at ($(L1.south)!0.5!(L2.north)$);
\coordinate (mid2) at ($(L2.south)!0.5!(L3.north)$);

% routed left of the lane frames so they never cross a lane label
\coordinate (offL) at ($(L2.west)-(3.5mm,0)$);

\draw[hand] (a5.south) -- (a5.south|-mid1) -- (offL|-mid1) -- (offL|-b1) -- (b1.west);
\node[above=0.2mm, font=\scriptsize, text=vizink]
  at ($(a5.south|-mid1)!0.5!(offL|-mid1)$) {\texttt{batch\_ids.json} via Actions cache};

\draw[hand] (b5.south) -- (b5.south|-mid2) -- (offL|-mid2) -- (offL|-c1) -- (c1.west);
\node[above=0.2mm, font=\scriptsize, text=vizink]
  at ($(b5.south|-mid2)!0.5!(offL|-mid2)$) {traces and scores via Langfuse API};

\end{tikzpicture}

    \end{adjustbox}
    \caption{Scheduled GitHub Actions workflows. Each lane is one workflow; the two blue handoffs are the only state shared between them, which is what allows the asynchronous batch APIs to be polled in a later job.}
    \label{fig:ga-weekly-pipelines}
\end{figure}

\subsection{Workflow~1: Submit Batches}
\label{sec:batch-process}

Langfuse holds \(\mathcal{D}_N\) and the frozen \(\mathcal{D}_n\) with stable \texttt{custom\_id} join keys and the boolean label \texttt{campaign\_relevant} (Chapter~\ref{ch:data}). Weekly runs submit only \(\mathcal{D}_n\).

Rows normalize to \(\{\texttt{custom\_id},\;\texttt{body},\;\texttt{campaign\_relevant}\}\), with \texttt{body} built from \texttt{config/request\_template.json} plus the multimodal \texttt{input}. The same template is reused every week so temporal differences are not request-construction artefacts. Orchestrator \texttt{run\_eval.py} loads the panel, ensures Gemini image uploads when needed, and dispatches each model to \texttt{batch\_*.py}. Adapters serialize to the provider batch format, upload JSONL, submit the job, and write the run manifest \texttt{data/batch\_ids.json} (\texttt{run\_id}, \texttt{subset\_size}, per-model batch IDs) into the Actions cache for retrieval.

The prediction task is identical across providers: JSON boolean \texttt{campaign\_relevant} under a schema-constrained output. Adapter differences are only transport-level (Table~\ref{tab:provider-adapters}). Claude sanitizes \texttt{custom\_id} on the wire and remaps on download. Gemini replaces image URLs with Files API URIs via \texttt{upload\_gemini\_images.py} and \texttt{data/gemini\_image\_cache.json} (TTL \(\approx\) 47~h). Across providers the output schema is \(\{\texttt{campaign\_relevant},\; \texttt{relevancy\_reasoning}\}\).

\begin{table}[htbp]
\centering
\footnotesize
\caption{Adapter-layer mapping for one shared task and one shared request body: what changes per provider and why.}
\label{tab:provider-adapters}
\setlength{\tabcolsep}{5pt}
\renewcommand{\arraystretch}{1.25}
\begin{tabular}{|p{2.0cm}|p{4.2cm}|p{8.1cm}|}
\hline
Provider & Schema / response constraint & Provider-specific preprocessing \\
\hline
OpenAI & \texttt{text.format} with JSON schema & Reuse canonical multimodal \texttt{input}; keep image URLs as-is \\
\hline
Gemini & \texttt{responseSchema} with JSON MIME type & Convert URLs to Files API \texttt{fileUri}; use upload cache (TTL \(\approx\) 47 h); skip rows with missing media \\
\hline
Claude & Anthropic \texttt{output\_config} schema & Sanitize \texttt{custom\_id} for submission; remap IDs on result download \\
\hline
\end{tabular}
\end{table}

\textbf{Output.} The submission workflow produces \texttt{data/batch\_ids.json} and the submitted JSONL files under \texttt{batches/}.

\subsection{Workflow~2: Retrieve and Ingest}
\label{sec:metrics-dashboard}

\texttt{poll\_and\_ingest.py} restores the manifest, polls until completion, and downloads JSONL into \texttt{data/results/}. Then \texttt{ingest.py} normalizes each response, joins \texttt{campaign\_relevant} to ground truth via \texttt{custom\_id}, maps token usage to \texttt{cost\_usd} via \texttt{config/models.yaml}, and writes Langfuse traces with scores.

Table~\ref{tab:batch-pricing-examples} lists batch prices for the monitored roster; \texttt{models.yaml} also carries unused candidates for future adds.

\begin{table}[htbp]
    \centering
    \caption{Batch-API prices per one million tokens (USD) for the monitored models, as of 16~March~2026. Claude batch prices are half the corresponding standard prices.}
    \label{tab:batch-pricing-examples}
    \begin{tabular}{lrr}
        \toprule
        Model & Input & Output \\
        \midrule
        \texttt{gpt-5.4-2026-03-05}      & 1.25  & 7.50  \\
        \texttt{gpt-5.2}                 & 0.875 & 7.00  \\
        \texttt{gpt-5.1}                 & 0.625 & 5.00  \\
        \texttt{gpt-5}                   & 0.625 & 5.00  \\
        \texttt{gpt-5-mini}              & 0.125 & 1.00  \\
        \texttt{gpt-5-nano}              & 0.025 & 0.20  \\
        \texttt{gpt-4.1-mini}            & 0.20  & 0.80  \\
        \texttt{gpt-4.1-nano}            & 0.05  & 0.20  \\
        \midrule
        \texttt{gemini-2.5-flash}        & 0.15  & 1.25  \\
        \midrule
        \texttt{claude-sonnet-4-6}       & 1.78 & 8.92  \\
        \texttt{claude-sonnet-4-5}       & 1.78 & 8.92  \\
        \texttt{claude-haiku-4-5}        & 0.59 & 2.97  \\
        \bottomrule
    \end{tabular}
\end{table}

Trace scores used later: \texttt{accuracy} (\(1/0\) match), \texttt{error\_type} (TP/FP/TN/FN), and \texttt{cost\_usd}. Precision, recall, and F1 follow Section~\ref{sec:classification-metrics}.

\subsection{Workflow~3: Dashboard Sync}

\texttt{sync-streamlit-dashboard.yml} exports traces/scores via \texttt{export\_langfuse\_csv.py} into committed CSVs in \texttt{gocomo/llm-batch-eval}, then mirrors \texttt{streamlit\_app/} to \texttt{GeorgiiBurdinComo/llm-eval-dashboard}. Tabs: Overview (accuracy/cost over time), McNemar (contingency tables and exact \(p\)), Utility (\(E_t[\mathcal{C}](m)\) with \(R=C_{\mathrm{FP}}/C_{\mathrm{FN}}\)).

\subsection{Artifacts and Cross-Step Data Flow}
\label{sec:artifacts-data-flow}

Workflows share only files: \texttt{data/batch\_ids.json}, \texttt{batches/}, \texttt{data/results/}, Langfuse, then CSV snapshots. Any stage can be replayed from the previous artifact without repeating provider calls.

\FloatBarrier

\section{Experimental Setup}
\label{sec:experimental-design}

Datasets and roster are those of Chapter~\ref{ch:data} and Table~\ref{tab:monitored-models}. This chapter reports the frozen-panel baseline, longitudinal McNemar evidence, and the corrective-action comparison panel. Figure~\ref{fig:run-decision-logic} is the operational loop; prompt optimisation is shown as a branch but only monitoring and comparison are evaluated here. The realized degradation episode is Chapter~\ref{chap:degradation-case-study}.

\FloatBarrier
\begin{figure}[H]
    \centering
    \begin{adjustbox}{max width=0.92\textwidth}
    \begin{tikzpicture}[
    >=Latex,
    node distance=9mm and 14mm
    ]
    
    \node[figbox] (baseline) {Baseline\\Run\\Metrics};
    \node[figbox, right=of baseline] (newrun) {New\\Run\\Metrics};
    
    \node[figbox, below=11mm of $(baseline)!0.5!(newrun)$] (drift) {DriftTest\\McNemar on fixed\\300-example subset};
    
    \node[figbox, left=24mm of drift] (deploy) {Deploy\\Current\\Setup};
    \node[figbox, right=20mm of drift] (reeval) {Escalation:\\Optional Full\\Re-evaluation\\$ \mathcal{D}_N $};
    \node[figbox, right=16mm of reeval] (metriccost) {Compare\\metrics + cost};
    
    \node[figdiamond, below=9mm of metriccost] (decision) {Choose\\corrective\\action};
    
    \node[figbox, below left=7mm and 12mm of decision] (switch) {Switch\\Model};
    \node[figbox, below right=7mm and 12mm of decision] (gepa) {Run\\GEPA};
    
    \draw[figarrow] (baseline) -- (drift);
    \draw[figarrow] (newrun) -- (drift);
    
    \draw[figarrow] (drift) -- node[midway, above, font=\scriptsize]{no drift} (deploy);
    \draw[figarrow] (drift) -- node[midway, above, font=\scriptsize]{drift + escalation} (reeval);
    
    \draw[figarrow] (reeval) -- (metriccost);
    \draw[figarrow] (metriccost) -- (decision);
    
    \draw[figarrow] (decision) -- (switch);
    \draw[figarrow] (decision) -- (gepa);
    
    \coordinate (loopabove) at ($(reeval.north)+(0,1.0)$);
    \draw[figarrow, rounded corners=2.5pt]
      (gepa.east) -- ++(0.9,0) |- (loopabove) -- (reeval.north);
    
    \end{tikzpicture}
    \end{adjustbox}
    \caption{Decision logic after a monitoring run: McNemar drift on the fixed \(n=300\) panel, optional full-benchmark re-evaluation on \(\mathcal{D}_N\), then metric/cost comparison and corrective action (model switch or GEPA loop).}
    \label{fig:run-decision-logic}
\end{figure}
\FloatBarrier

\section{Results}

\subsection{Baseline and Instantiated Monitoring Subset}
\label{sec:baseline-subset}

A one-time \(\mathcal{D}_N\) run seeds the disagreement construction of Chapter~\ref{ch:data}; the resulting \(n=300\) panel is frozen in Langfuse. Table~\ref{tab:phase1-baseline-snapshot} is a descriptive baseline excerpt on \(\mathcal{D}_n\) (metrics export, last score day 2026-04-13).

\begin{table}[htbp]
    \centering
    \caption{Baseline excerpt: classification metrics on the frozen monitoring subset (metrics export, last score day 2026-04-13).}
    \label{tab:phase1-baseline-snapshot}
    \begin{tabular}{lccccc}
        \toprule
        Model & $n$ & Accuracy & Precision & Recall & F1 \\
        \midrule
        \texttt{gpt-4.1-mini} & 300 & 0.8033 & 0.9948 & 0.7680 & 0.8668 \\
        \texttt{gpt-4.1-nano} & 300 & 0.7000 & 0.9444 & 0.6800 & 0.7907 \\
        \bottomrule
    \end{tabular}
\end{table}

The baseline fixes the initial error profile against which later runs are compared; fallback choice depends on which error type is costly under Chapter~\ref{ch:cost-deployment-utility}.

\subsection{Realized discordance}
\label{sec:realized-discordance}

The panel size in Section~\ref{sec:monitoring-sample-size} was derived from a pilot discordance rate of $\hat\rho=0.114$, measured by running $\mathcal{D}_N$ twice two days apart. Table~\ref{tab:discordance} reports the discordance actually observed between each model's baseline run and its later runs.

\begin{table}[htbp]
    \centering
    \footnotesize
    \caption{Discordance rate $(b+c)/n$ between each model's onset baseline and later
    scheduled panel runs. Pooled median 0.027.}
    \label{tab:discordance}
    \begin{tabular}{lcccc}
    \toprule
    Model & Comparisons & Median & Min & Max \\
    \midrule
    \texttt{gpt-4.1-nano} & 16 & 0.152 & 0.124 & 0.187 \\
    \texttt{gpt-5-nano} & 15 & 0.124 & 0.094 & 0.149 \\
    \texttt{gpt-5-mini} & 15 & 0.080 & 0.070 & 0.104 \\
    \texttt{gpt-4.1-mini} & 16 & 0.062 & 0.040 & 0.206 \\
    \texttt{claude-haiku-4-5} & 14 & 0.034 & 0.015 & 0.049 \\
    \texttt{claude-sonnet-4-5} & 14 & 0.030 & 0.023 & 0.053 \\
    \texttt{gpt-5.1} & 15 & 0.020 & 0.010 & 0.034 \\
    \texttt{gemini-2.5-flash} & 12 & 0.017 & 0.010 & 0.023 \\
    \texttt{claude-sonnet-4-6} & 14 & 0.015 & 0.008 & 0.023 \\
    \texttt{gpt-5.4} & 13 & 0.013 & 0.007 & 0.017 \\
    \texttt{gpt-5.2} & 14 & 0.010 & 0.003 & 0.013 \\
    \texttt{gpt-5} & 15 & 0.007 & 0.003 & 0.010 \\
    \bottomrule
    \end{tabular}
\end{table}


The pooled median is $0.033$, well below the pilot value. Two things drive the gap. First, the pilot measured one model; the roster spans a wide stability range, from \texttt{gpt-4.1-nano} at $0.156$ down to \texttt{gpt-5} at $0.007$. Second, the pilot model was the production model of the time, which sat in the churn-heavy part of that range.

The consequence is a power loss concentrated in the stable models. The planning equation~\eqref{eq:planning-discordant-m} needs $m\approx 31$ discordant pairs to detect $\delta=0.05$ at $80\%$ power. At $\rho=0.007$, a 300-item panel yields $m\approx 2$. For \texttt{gpt-5} and \texttt{gpt-5.2} the panel is therefore too small to detect a five-point drop, and their flat monitoring series should be read as uninformative rather than as evidence of stability. For the models at $\rho\ge 0.10$ the panel delivers $m\approx 30$ or more and the sizing assumption holds.

The panel was sized once, for one model, and then frozen. Per-model sizing --- a larger panel for low-churn models --- would fix this, at proportionally higher cost per run.

\subsection{Longitudinal Monitoring Results}
\label{sec:drift-results}

This subsection checks whether repeated evaluation on the fixed monitoring subset shows meaningful temporal changes in model behaviour. Since the panel is fixed, changes in the paired discordance counts \((b,c)\) reflect changed predictions rather than a changed sample.

\begin{table}[htbp]
    \centering
    \caption{McNemar table for \texttt{gpt-4.1-mini} between 6~April~2026 and 13~April~2026 on \(\mathcal{D}_n\) ($n=299$).}
    \label{tab:mcnemar_gpt41mini}
    \begin{tabular}{lcc}
        \toprule
        & 2026-04-13 correct & 2026-04-13 wrong \\
        \midrule
        2026-04-06 correct & 230 & 10 \\
        2026-04-06 wrong   & 10  & 49 \\
        \bottomrule
    \end{tabular}
\end{table}

\begin{table}[htbp]
    \centering
    \caption{McNemar table for \texttt{gpt-4.1-nano} between 6~April~2026 and 13~April~2026 on \(\mathcal{D}_n\).}
    \label{tab:mcnemar_gpt41nano}
    \begin{tabular}{lcc}
        \toprule
        & 2026-04-13 correct & 2026-04-13 wrong \\
        \midrule
        2026-04-06 correct & 184 & 22 \\
        2026-04-06 wrong   & 26  & 68 \\
        \bottomrule
    \end{tabular}
\end{table}

Two recent run pairs show the range of outcomes. For \texttt{gpt-4.1-mini}, the comparison from 6~April~2026 to 13~April~2026 gives the contingency table in Table~\ref{tab:mcnemar_gpt41mini}. The discordant counts are \((b,c)=(10,10)\), so gains and losses balance exactly and there is no directional change.

For \texttt{gpt-4.1-nano}, the corresponding comparison in Table~\ref{tab:mcnemar_gpt41nano} gives \((b,c)=(22,26)\): 48 instances change status, with 22 regressions and 26 recoveries. The one-sided McNemar test does not reach \(\alpha=0.05\). Thus, despite noticeable churn, the net direction is not degradation under the alert rule from Chapter~\ref{ch:drift}.

These two examples show two different situations on the fixed panel: balanced movement with no directional evidence (\texttt{gpt-4.1-mini}) and high-churn but non-degrading movement (\texttt{gpt-4.1-nano}). In both cases, the correct operational response is continued monitoring together with preparation of fallback comparisons on the metric and cost panel.

\subsection{Alert counts over the full monitoring period}
\label{sec:method-comparison}

The monitoring series covers 8~March to 27~July~2026 and yields 265 baseline-to-current comparisons over 19 models. Twelve model-days are full-benchmark runs of about $1{,}190$ items rather than panel runs; these are escalation and prompt-transfer evaluations, not part of the weekly loop, and they are excluded from the drift series. Mixing them in would compare a panel run against a benchmark run and attribute the difference to the model.

\begin{table}[htbp]
    \centering
    \caption{Alert counts over 173 baseline-to-later comparisons
    (12 models, scheduled panel runs, 15~March--6~July~2026).
    Holm correction is applied once to the full primary family.}
    \label{tab:alert-counts}
    \begin{tabular}{lr}
    \toprule
    Rule & Alerts \\
    \midrule
    $b>c$, $p_{\mathrm{exact}}<0.05$ (uncorrected) & 3 \\
    $b>c$, Holm-adjusted $p<0.05$ & 1 \\
    Holm and $\Delta A_{\mathrm{paired}}\le-3$pp (operational) & 1 \\
    \midrule
    $\Delta A_{\mathrm{paired}}\le-3$pp alone & 2 \\
    $\Delta A_{\mathrm{paired}}\le-5$pp alone & 1 \\
    \bottomrule
    \end{tabular}
\end{table}


Table~\ref{tab:alert-counts} gives the alert counts. At the nominal $\alpha=0.05$ the rule fires 13 times. The Bonferroni correction of~\eqref{eq:bonferroni} reduces this to 3. Adding a three-point relevance floor removes nothing further: all three surviving comparisons are drops of $4.7$pp or more, so on this data the floor is inert and the correction is doing the filtering.

The three surviving comparisons are not three episodes. Two of them are \texttt{gpt-4.1-mini} against its 8~March baseline on 22 and 23~March, which is one event observed on two consecutive days; Chapter~\ref{chap:degradation-case-study} treats it in full. The third is \texttt{gpt-5-mini} on 6~April, at $-4.7$pp with $p=1.3\times10^{-3}$, and it needs a different reading.

\paragraph{The \texttt{gpt-5-mini} alert is baseline noise.}
\texttt{gpt-5-mini} scores $0.940$ on the 8~March baseline and then sits between $0.893$ and $0.933$ for the following twenty-two runs, with no trend. The baseline is the highest single value in the series. Every later comparison therefore inherits a small negative offset, and nine of the twenty-two clear $p<0.05$ before correction. The 6~April run is simply the lowest of them. Reading this as a provider-side change would be wrong: nothing steps, and the post-baseline level is flat.

This is the cost of the fixed-baseline design noted in Section~\ref{sec:baseline-multiplicity}. Anchoring on one batch makes the entire series inherit that batch's sampling error, and a baseline that happens to land high manufactures a persistent apparent decline. Bonferroni suppresses eight of the nine, but not the largest, so multiplicity control alone does not remove the problem. Two mitigations are available and neither was in place here: pool the first two runs into the baseline, or require an alert to repeat before it is acted on. On this data the second rule would have retained the \texttt{gpt-4.1-mini} event, which fired on two consecutive days, and dropped the \texttt{gpt-5-mini} alert, which fired once.

By simple thresholds, a $-3$pp rule fires 15 times and a $-5$pp rule twice. The $-5$pp rule selects exactly the two \texttt{gpt-4.1-mini} comparisons --- on this data the crudest rule available recovers the one real event and nothing else. The paired test earns its place not through a lower alert count but by reporting $b$ and $c$ separately, which is what makes the \texttt{gpt-5-mini} case diagnosable at all: a threshold rule reports only that accuracy is down.

\subsection{Candidate Comparison for Corrective Action}
\label{sec:cost-aware-comparison}

After escalation, candidates are compared on accuracy, precision, recall, F1, and mean cost. Table~\ref{tab:model-metrics-snapshot} is the \(\mathcal{D}_n\) profile (confusion counts plus derived metrics; \texttt{cost\_usd} from Section~\ref{sec:metrics-dashboard}). Final expected-cost ranking uses \(\mathcal{D}_N\) under Chapter~\ref{ch:cost-deployment-utility} (Section~\ref{sec:selection-bias}).

\begin{table}[htbp]
    \centering
    \footnotesize
    \caption{Exported snapshot of confusion counts and classification metrics per candidate (evaluation on \(\mathcal{D}_n\)).}
    \label{tab:model-metrics-snapshot}
    \begin{adjustbox}{max width=\linewidth}
    % Auto-generated — last score day 2026-03-28
% % Auto-generated — last score day 2026-03-28
% % Auto-generated — last score day 2026-03-28
% \input{metrics_tabular.tex} inside your table environment
\begin{tabular}{lrrrrrrrrrrrr}
\toprule
model & n & tp & fp & tn & fn & accuracy & precision & recall & f1 & specificity & balanced\_accuracy & pr\_auc \\
\midrule
gpt-4.1-mini & 299 & 190 & 1 & 49 & 59 & 0.7993 & 0.9948 & 0.7631 & 0.8636 & 0.9800 & 0.8715 & 0.9776 \\
gpt-4.1-nano & 300 & 178 & 8 & 42 & 72 & 0.7333 & 0.9570 & 0.7120 & 0.8165 & 0.8400 & 0.7760 & 0.9545 \\
\bottomrule
\end{tabular}

 inside your table environment
\begin{tabular}{lrrrrrrrrrrrr}
\toprule
model & n & tp & fp & tn & fn & accuracy & precision & recall & f1 & specificity & balanced\_accuracy & pr\_auc \\
\midrule
gpt-4.1-mini & 299 & 190 & 1 & 49 & 59 & 0.7993 & 0.9948 & 0.7631 & 0.8636 & 0.9800 & 0.8715 & 0.9776 \\
gpt-4.1-nano & 300 & 178 & 8 & 42 & 72 & 0.7333 & 0.9570 & 0.7120 & 0.8165 & 0.8400 & 0.7760 & 0.9545 \\
\bottomrule
\end{tabular}

 inside your table environment
\begin{tabular}{lrrrrrrrrrrrr}
\toprule
model & n & tp & fp & tn & fn & accuracy & precision & recall & f1 & specificity & balanced\_accuracy & pr\_auc \\
\midrule
gpt-4.1-mini & 299 & 190 & 1 & 49 & 59 & 0.7993 & 0.9948 & 0.7631 & 0.8636 & 0.9800 & 0.8715 & 0.9776 \\
gpt-4.1-nano & 300 & 178 & 8 & 42 & 72 & 0.7333 & 0.9570 & 0.7120 & 0.8165 & 0.8400 & 0.7760 & 0.9545 \\
\bottomrule
\end{tabular}


    \end{adjustbox}
\end{table}

Read confusion counts first, then derived metrics, then \(E_t[\mathcal{C}]\) on \(\mathcal{D}_N\). Similar accuracy can hide very different FP/FN mixes; which model wins depends on campaign costs via \(R=C_{\mathrm{FP}}/C_{\mathrm{FN}}\).
\section{Limitations}
\label{sec:error-analysis}

\begin{itemize}
    \item Provider state is unobserved; drift is inferred only from output changes on the fixed panel.
    \item Cost depends on provider pricing and tokenization, so monetary comparisons may change even when the confusion structure stays the same.
    \item Every comparison reuses one baseline run, so sampling error in that run propagates to the whole series for that model (Section~\ref{sec:method-comparison}).
    \item The panel was sized from a single pilot discordance estimate. For the most stable models the realized discordance leaves the panel underpowered at the design effect size (Section~\ref{sec:realized-discordance}).
\end{itemize}

The degradation episode itself is Chapter~\ref{chap:degradation-case-study}.
